\subsection{伽罗瓦理论的基本定理}

定理 4.——设 N 是 K 的一个伽罗瓦扩张，\(\Gamma\) 是其伽罗瓦群。设 \(\mathcal{K}\) 为 N 的子扩张组成的集，\(\mathcal{G}\) 为 \(\Gamma\) 的闭子群组成的集。对每个子群 \(\Delta \in \mathcal{G}\) ，我们用 \(k(\Delta)\) 表示 \(\Delta\) 的不变量域；对每个子域 E \(\in \mathcal{K}\) ，我们用 g(E) 表示 N 的 E-自同构组成的群。则 \(\Delta \mapsto k(\Delta)\) 是 \(\mathcal{G}\) 到 \(\mathcal{K}\) 上的一个双射，而 E \(\mapsto\) g(E) 是其逆双射。

\begin{enumerate}
\def\labelenumi{\Alph{enumi})}
\tightlist
\item
  关系 \(E = k(g(E))\) （对 \(E \in \mathcal{K}\) ）是下述更精确的引理的推论：
\end{enumerate}

引理 1.——设 E 是 N 的一个子扩张。则 N 是 E 的一个伽罗瓦扩张，且 Gal(N/E) 是 Gal(N/K) 的一个闭子群（带诱导拓扑）。

设 \(x \in \mathbb{N}\) ；极小多项式 \(f\) （\(x\) 的，\(\mathbf{E}\) 上的）在 \(\mathbf{E}[\mathbf{X}]\) 中整除极小多项式 \(g\) （\(x\) 的，\(\mathbf{K}\) 上的）（V，第 17 页，推论 2）。由于 \(\mathbf{N}\) 在 \(\mathbf{K}\) 上是伽罗瓦的，多项式 \(g\) 在 \(\mathbf{N}[\mathbf{X}]\) 中是互异的一次因子的乘积；故 \(f\) 亦然，从而 \(\mathbf{N}\) 在 \(\mathbf{E}\) 上是伽罗瓦的。

设 \(\Gamma\) 是 N 在 K 上的伽罗瓦群，\(\Delta\) 是 N 在 E 上的伽罗瓦群。按定义，\(\Delta\) 是 \(\Gamma\) 的由所有这样的 \(\sigma\) 组成的子群：\(\sigma(x) = x\) （对所有 \(x \in E\) ）。而对每个 \(x \in E\) ，映射 \(\sigma \mapsto \sigma(x)\) （\(\Gamma\) 到离散空间 N 中的）是连续的，故 \(\Delta\) 在 \(\Gamma\) 中是闭的。设 \(\sigma \in \Gamma\) ；对 N 中的 \(x_1, \ldots, x_n\) ，令 U \((x_1, \ldots, x_n)\) 为所有这样的 \(\tau \in \Gamma\) 组成的集：\(\tau(x_i) = \sigma(x_i)\) （对 \(1 \leqslant i \leqslant n\) ）；令

\[
\mathrm{V} \left(x _ {1}, \dots , x _ {n}\right) = \mathrm{U} \left(x _ {1}, \dots , x _ {n}\right) \cap \Delta .
\]

则集族 \(\mathrm{U}(x_{1},\ldots,x_{n})\) （相应地，\(\mathrm{V}(x_{1},\ldots,x_{n})\) ）是 \(\sigma\) 在 \(\Gamma\) 中（相应地，在 \(\Delta\) 中）的邻域基。故 \(\Delta\) 上的拓扑是由 \(\Gamma\) 诱导的拓扑。

\begin{enumerate}
\def\labelenumi{\Alph{enumi})}
\setcounter{enumi}{1}
\tightlist
\item
  关系 \(\Delta = g(k(\Delta))\) （对 \(\Delta \in \mathcal{G}\) ）是下述更精确的引理的推论：
\end{enumerate}

引理 2.——设 \(\Delta\) 是 \(\Gamma\) 的一个子群。设 \(E\) 为 \(\Delta\) 的不变量域；则 \(N\) 在 \(E\) 上的伽罗瓦群是 \(\Delta\) 在 \(\Gamma\) 中的闭包。

N 在 E 上的伽罗瓦群在 \(\Gamma\) 中是闭的（引理 1）且包含 \(\Delta\) ，故它包含闭包 \(\bar{\Delta}\) （\(\Delta\) 的）。设 \(\sigma\) 是 N 的一个 E-自同构，设 \(x_{1}, \ldots, x_{n}\) 在 N 中。由于 N 在 E 上是伽罗瓦的（引理 1），存在（V，第 57 页，命题 2）N 的一个子扩张 \(N_{0}\) ，它在 E 上伽罗瓦、次数有限，且包含 \(x_{1}, \ldots, x_{n}\) 。设 \(\Delta_{0}\) 为子群 \(\Delta\) （\(\operatorname{Gal}(N/E)\) 的）在 \(\operatorname{Gal}(N/E)\) 到 \(\operatorname{Gal}(N_{0}/E)\) 中的限制同态下的像。由于 \([N_{0}:E]\) 有限，戴德金定理（V，第 27 页，推论 2）表明 \(\operatorname{Gal}(N_{0}/E)\) 有限。故 \(\Delta_{0}\) 有限，又因 E 是 \(\Delta_{0}\) 的不变量域，我们有 \(\Delta_{0}=\)

Gal(N \(_0\) /E)（V，第 66 页，定理 3）。特别地，\(\Delta_0\) 包含 \(\sigma\) 在 N \(_0\) 上的限制。因此存在 \(\tau \in \Delta\) 使 \(\sigma\) 与 \(\tau\) 在 N \(_0\) 上有相同的限制，从而 \(\sigma(x_1) = \tau(x_1)\) , \ldots, \(\sigma(x_n) = \tau(x_n)\) 。由此得出 \(\sigma\) 是 \(\Delta\) 在 \(\Gamma\) 中的一个极限点，从而 Gal(N/E) \(\subset \bar{\Delta}\) .

推论 1.——设 E 与 E' 是 N 的包含 K 的两个子域；则 E ⊂ E' 当且仅当 g(E) ⊃ g(E')。若 Δ 与 Δ' 是 Γ 的两个闭子群，则 Δ ⊂ Δ' 当且仅当 k(Δ) ⊃ k(Δ')。

因为这两个互逆的双射 \(E \mapsto g(E)\) 与 \(\Delta \mapsto k(\Delta)\) 是反转包含关系的。

推论 2.——设 \((\mathrm{E}_i)_{i\in I}\) 是 N 的包含 K 的子域的一个族；令 \(\mathbf{L} = \bigcap_{i\in I}\mathbf{E}_i\) ，\(\mathbf{M} = \mathbf{K}\left(\bigcup_{i\in I}\mathbf{E}_i\right)\) 。则 \(g(L)\) 是 \(\Gamma\) 的最小闭子群，

它包含 \(\bigcup_{i\in I}g(E_i)\) ，且我们有 \(g(M) = \bigcap_{i\in I}g(E_i)\) 。

第一个断言由推论 1 得出，第二个是直接的。

推论 3.——对 \(i=1\) , 2，设 \(\mathrm{E}_{i}\) 是 N 的包含 K 的一个子域，并设 \(\Delta_{i}=g(\mathrm{E}_{i})\) 。对任意 \(\sigma\in\Gamma\) ，关系 \(\sigma(\mathrm{E}_{1})=\mathrm{E}_{2}\) 与 \(\sigma\Delta_{1}\sigma^{-1}=\Delta_{2}\) 等价。

因为我们有 \(\tau \in g(\sigma(E_1))\) 当且仅当 \(\tau\sigma(x) = \sigma(x)\) ，即 \(\sigma^{-1}\tau\sigma(x) = x\) ，对所有 \(x \in E_1\) ；这等于说 \(\sigma^{-1}\tau\sigma \in \Delta_1\) ，从而 \(g(\sigma(E_1)) = \sigma\Delta_1\sigma^{-1}\) 。

推论 4.——设 E 是 N 的包含 K 的一个子域，并设 \(\Delta = g(E)\) 。E 在 K 上是伽罗瓦的，当且仅当 \(\Delta\) 是 \(\Gamma\) 的一个正规子群。当此条件成立时，\(\Gamma\) 到 Gal(E/K) 中的限制同态通过过渡到商群定义了 \(\Gamma/\Delta\) 到 Gal(E/K) 上的一个拓扑群同构。

由于 N 在 K 上可分，E 亦然（V，第 36 页，命题 1）。因此 E 在 K 上是伽罗瓦的，当且仅当它在 K 上是拟伽罗瓦的；这也意味着 \(\sigma(E) = E\) 对 N 的每个 K-自同构 \(\sigma\) 成立（V，第 52 页，命题 1 及第 54 页，命题 3）。由推论 3，这等价于 \(\sigma\Delta\sigma^{-1} = \Delta\) （对所有 \(\sigma \in \Gamma\) ）。

限制同态 \(\varphi: \operatorname{Gal}(N/K) \to \operatorname{Gal}(E/K)\) 是连续且满的（V，第 60 页，命题 3），其核显然等于 \(\Delta = \operatorname{Gal}(N/E)\) 。由于 \(\Gamma\) 紧，\(\Gamma/\Delta\) 到 \(\operatorname{Gal}(E/K)\) 上的、由 \(\varphi\) 通过过渡到商群得到的同态是拓扑群的同构（一般拓扑学，I，第 87 页，推论 2）。

推论 5.——设 E 是 N 的包含 K 的一个子域。E 在 K 上次数有限，当且仅当 \(g(E)\) 在 \(\Gamma\) 中是开的。当此条件成立时，指标 \((\Gamma:g(E))\) 有限且等于 \([E:K]\) 。

\(g(E)\) 是开的，当且仅当存在 N 的一个子扩张 F，在 K 上次数有限，使得按 V 第 60 页的记号，\(g(E)\) 包含 \(U_{F}(Id_{N}) = g(F)\) 。关系 \(g(E) \supset g(F)\) 由推论 1（V，第 68 页）等价于 \(E \subset F\) ，由此得推论 5 的第一个断言。

设 \([E:K]\) 有限。设 \(\Omega\) 为 K 的一个把 N 作为子扩张包含在内的代数闭包（V，第 23 页，定理 2），设 H 为 E 到 \(\Omega\) 中的 K-同态组成的集。H 的每个元素由 \(\Omega\) 的一个 K-自同构诱导（V，第 52 页，命题 1），且由于 N 在 K 上是拟伽罗瓦的，映射 \(\sigma\mapsto\sigma|E\) （\(\Gamma\) 到 H 上的）是满的。\(\sigma\) 与 \(\sigma'\) （\(\Gamma\) 中的）在 E 上有相同的限制，当且仅当 \(\sigma^{-1}\sigma'\in g(E)\) ，由此 Card \(H=(\Gamma:g(E))\) 。最后，由于 E 是 K 上的一个平展代数，我们有 Card \(H=[E:K]\) （V，第 32 页，命题 4），故最终我们有 \((\Gamma:g(E))=[E:K]\) 。

推论 6.——对 \(i=1,2\) ，设 \(\mathbf{E}_{i}\) 是 N 的一个子扩张，\(\Gamma_{i}\) 是 N 在 \(\mathbf{E}_{i}\) 上的伽罗瓦群。下列条件等价：

\begin{enumerate}
\def\labelenumi{\alph{enumi})}
\item
  群 \(\Gamma\) 是子群 \(\Gamma_1\) 与 \(\Gamma_2\) 的直积。
\item
  扩张 \(\mathbf{E}_1\) 与 \(\mathbf{E}_2\) 在 \(\mathbf{K}\) 上是伽罗瓦的，我们有 \(\mathbf{E}_1 \cap \mathbf{E}_2 = \mathbf{K}\) ，且
\end{enumerate}

\[
\mathbf {K} \left(\mathrm{E} _ {1} \cup \mathrm{E} _ {2}\right) = \mathbf {N}.
\]

\(\Gamma\) 是子群 \(\Gamma_{1}\) 与 \(\Gamma_{2}\) 的直积，当且仅当下列条件成立（I，第 48 页，命题 15）：

\begin{enumerate}
\def\labelenumi{(\roman{enumi})}
\item
  子群 \(\Gamma_{1}\) 与 \(\Gamma_{2}\) 在 \(\Gamma\) 中是正规的；
\item
  \(\Gamma_1 \cap \Gamma_2 = \{\varepsilon\}\) ，其中 \(\varepsilon\) 是 \(\Gamma\) 的单位元；
\item
  \(\Gamma = \Gamma_{1} \cdot \Gamma_{2}\) 。
\end{enumerate}

今 (i) 意味着 \(E_{1}\) 与 \(E_{2}\) 在 K 上是伽罗瓦的（推论 4）。由推论 2，条件 (ii) 等价于 \(\mathbf{N} = \mathbf{K}(\mathbf{E}_{1} \cup \mathbf{E}_{2})\) ；最后，若 (i) 与 (ii) 成立，则 \(\Gamma_{1}\Gamma_{2}\) 是 \(\Gamma\) 的包含 \(\Gamma_{1} \cup \Gamma_{2}\) 的最小子群；它是闭的，因为 \(\Gamma_{1}\) 与 \(\Gamma_{2}\) 紧，且映射 \((\sigma, \tau) \mapsto \sigma\tau\) （\(\Gamma_{1} \times \Gamma_{2}\) 到 \(\Gamma\) 中的）是连续的（一般拓扑学，I，第 63 页，推论 1）。于是推论 2 表明 (iii) 等价于 \(E_{1} \cap E_{2} = K\) ，这就证明了 a) 与 b) 的等价性。

注.——用推论 6 的记号，设条件 a) 与 b) 成立。限制同态 \(\varphi_{i}:\Gamma\to\mathrm{Gal}(\mathrm{E}_{i}/\mathrm{K})\) （对 \(i=1,2\) ）诱导出拓扑群同构

\[
\Psi_ {1}: \Gamma_ {2} \rightarrow \operatorname{Gal} \left(\mathrm{E} _ {1} / \mathrm{K}\right), \quad \Psi_ {2}: \Gamma_ {1} \rightarrow \operatorname{Gal} \left(\mathrm{E} _ {2} / \mathrm{K}\right).
\]

由 \(a\) 我们看到，映射 \(\sigma \mapsto (\varphi_1(\sigma), \varphi_2(\sigma))\) 是 \(\operatorname{Gal}(\mathbf{N}/\mathbf{K})\) 到 \(\operatorname{Gal}(\mathrm{E}_1/\mathrm{K}) \times \operatorname{Gal}(\mathrm{E}_2/\mathrm{K})\) 上的一个拓扑群同构。
% semantic-fragments: legacy-input-seam
\relax{}
